\section{Receiver-factored architectural presentations}\label{sec:receiver-foundation}

This section introduces the minimal object required by the paper. The represented update is assumed. The architecture consists of a marked decomposition of the successor into receiver obligations and a marked factorization of each receiver branch.

\subsection{Represented updates and an upstream barrier}

Fix a represented update
\[
F:X\to Y.
\]
The spaces $X$ and $Y$ are already the states on which the architecture acts. When an upstream representation map
\[
\rho:\Omega\twoheadrightarrow X
\]
is relevant, its fibers impose a prior barrier, but $\rho$ is not part of the architectural primitive studied below.

\begin{lemma}[Fiber factorization criterion]\label{lem:fiber-factorization}
Let $Q:X\twoheadrightarrow Z$ be surjective and let $B:X\to W$. There exists a unique $G:Z\to W$ such that $B=GQ$ if and only if
\[
Q(x)=Q(x')\Longrightarrow B(x)=B(x').
\]
\end{lemma}

\begin{proof}
If $B=GQ$, equality of interface values implies equality of branch values. Conversely, if $B$ is constant on every $Q$-fiber, define $G(z)=B(x)$ for any $x$ with $Q(x)=z$. The fiber condition makes $G$ well defined, and surjectivity gives uniqueness.
\end{proof}

\begin{corollary}[Upstream representation barrier]\label{cor:representation-barrier}
A target $f:\Omega\to W$ factors through $\rho$ if and only if it is constant on the fibers of $\rho$. No downstream architecture operating only on $X$ can recover distinctions erased by $\rho$.
\end{corollary}

The remainder of the paper is conditional on $X$ and begins after this upstream identification has already occurred.

\subsection{Marked receiver presentations}

\begin{definition}[Marked receiver presentation]\label{def:marked-receiver-presentation}
A marked receiver presentation of $Y$ consists of:
\begin{enumerate}[label=(\roman*)]
\item a finite receiver set $R$;
\item receiver value spaces $(W_j)_{j\in R}$;
\item a bijection
\[
\chi_Y:Y\xrightarrow{\cong}H_Y\subseteq\prod_{j\in R}W_j;
\]
\item distinguished receiver projections $\pi_j:H_Y\to W_j$.
\end{enumerate}
\end{definition}

The bijection makes this a complete presentation of the represented successor state rather than a further quotient of it. If a scientific question intentionally replaces $Y$ by a quotient $\bar Y$, the definition is applied to the induced update $\bar F:X\to\bar Y$. A noninjective map out of the original $Y$ would instead specify only a partial or quotient target profile and would not be a complete receiver presentation of $F$.

The presentation is marked because the projections are part of the architectural claim. A global recoding of $Y$ can preserve total-state information while destroying the identification of one coordinate as receiver $j$.

\begin{definition}[Receiver branch]\label{def:receiver-branch}
For receiver $j\in R$, define
\[
B_j:=\pi_j\chi_YF:X\to W_j.
\]
\end{definition}

\begin{proposition}[Full-access receiver collapse]\label{prop:full-access-collapse}
Every branch has the trivial factorization $B_j=B_j\Id_X$. Moreover, the family $(B_j)_{j\in R}$ is uniquely equivalent to the product-valued map
\[
B:=\langle B_j\rangle_{j\in R}:X\to\prod_{j\in R}W_j,
\]
whose image lies in $H_Y$. Because $\chi_Y$ is bijective, $F=\chi_Y^{-1}B$. Thus receiver labels with complete shared access impose no nontrivial access architecture beyond the marked successor presentation itself.
\end{proposition}

\begin{proof}
The first claim is immediate. The product map is defined by $B(x)=(B_j(x))_j$ and is uniquely determined by its projections. By construction, $B=\chi_YF$, so its image lies in $H_Y$ and applying $\chi_Y^{-1}$ recovers $F$.
\end{proof}

\subsection{The architectural object}

\begin{definition}[Receiver-factored architectural presentation]\label{def:receiver-factored-architecture}
Let $F:X\to Y$ have the marked receiver presentation of \Cref{def:marked-receiver-presentation}. A receiver-factored architectural presentation is
\[
\mathcal A
=
\left(
F,\chi_Y,
\{X\xrightarrow{Q_j}Z_j\xrightarrow{G_j}W_j\}_{j\in R}
\right),
\]
where each $Q_j:X\twoheadrightarrow Z_j$ is surjective and
\begin{equation}\label{eq:receiver-factorization}
\boxed{
B_j
=
\pi_j\chi_YF
=
G_jQ_j.
}
\end{equation}
The map $Q_j$ is the receiver interface and $G_j$ is the receiver-local continuation.
\end{definition}

Every field in this object is typed and used in \eqref{eq:receiver-factorization}. The definition does not add a catch-all space of messages, graphs, procedures, or presentations. Those appear only when a particular $Q_j$ or $G_j$ is refined.

\begin{proposition}[Receiver-interface re-presentation]\label{prop:interface-representation}
Let $B_j=G_jQ_j$ with $Q_j:X\twoheadrightarrow Z_j$, and let $\zeta_j:Z_j\to Z'_j$ be a bijection. Define
\[
Q'_j:=\zeta_jQ_j,
\qquad
G'_j:=G_j\zeta_j^{-1}.
\]
Then $B_j=G'_jQ'_j$, and $Q_j$ and $Q'_j$ induce the same partition of $X$. We treat them as coordinate re-presentations of the same receiver interface.
\end{proposition}

\begin{proof}
The branch identity follows by cancellation of $\zeta_j^{-1}\zeta_j$. Because $\zeta_j$ is injective,
\[
Q'_j(x)=Q'_j(x')
\quad\Longleftrightarrow\quad
Q_j(x)=Q_j(x'),
\]
so the two interfaces have exactly the same fibers.
\end{proof}

This convention concerns only bijective recoding of one receiver interface. Comparisons that also transport receiver systems, schedules, interventions, implementation witnesses, or tracked roles require additional marked structure beyond the present paper.

\begin{definition}[Receiver indistinguishability]\label{def:receiver-indistinguishability}
For receiver $j$, define
\[
x\sim_{Q_j}x'
\quad\Longleftrightarrow\quad
Q_j(x)=Q_j(x').
\]
\end{definition}

By \Cref{lem:fiber-factorization}, every continuation through $Q_j$ is constant on these equivalence classes. Receiver locality is therefore exact relative to the marked interface: it is a declaration of which predecessor distinctions can affect that branch through the specified continuation.

\subsection{The canonical answer quotient is not the architecture}

For one branch $B:X\to W$, define
\[
x\sim_Bx'
\quad\Longleftrightarrow\quad
B(x)=B(x')
\]
and let $q_B:X\twoheadrightarrow X/{\sim_B}$ be the quotient.

\begin{proposition}[Canonical exact branch quotient]\label{prop:canonical-branch-quotient}
The branch $B$ factors uniquely through $q_B$. If $Q:X\twoheadrightarrow Z$ is any exact receiver interface for $B$, then there is a unique $h:Z\to X/{\sim_B}$ such that
\[
q_B=hQ.
\]
Hence $q_B$ is the coarsest exact quotient sufficient for the selected branch.
\end{proposition}

\begin{proof}
Define $\bar B([x])=B(x)$. If $B=GQ$, define $h(Q(x))=[x]$. Exactness implies that $Q(x)=Q(x')$ only when $B(x)=B(x')$, so $h$ is well defined. Surjectivity gives uniqueness.
\end{proof}

The two extensional extremes are
\[
Q=\Id_X
\qquad\text{and}\qquad
Q=q_B.
\]
The first retains everything; the second retains only the answer-equivalence class. Neither identifies the intermediate organization actually used. A receiver interface is therefore marked independently of the branch it realizes.

\subsection{Interface refinement}

\begin{definition}[Interface refinement]\label{def:interface-refinement}
For surjective interfaces $Q:X\twoheadrightarrow Z$ and $Q':X\twoheadrightarrow Z'$, write
\[
Q\preceq Q'
\]
when there is a map $h:Z'\to Z$ with $Q=hQ'$. Then $Q'$ retains every distinction retained by $Q$ and possibly more.
\end{definition}

\begin{proposition}[Refinement preorder and branch monotonicity]\label{prop:refinement-preorder}
The relation $\preceq$ is reflexive and transitive. If $Q\preceq Q'$ and a branch factors through $Q$, then it also factors through $Q'$.
\end{proposition}

\begin{proof}
Reflexivity uses the identity; transitivity uses composition. If $Q=hQ'$ and $B=GQ$, then $B=(Gh)Q'$.
\end{proof}

\subsection{Locality is presentation-relative}

\begin{proposition}[Global conjugacy can destroy coordinate locality]\label{prop:locality-conjugacy}
Let $Y=\R^2$ with marked receiver coordinates and define
\[
F(x,y)=(x^2,y),
\qquad
T(x,y)=(x+y,x-y).
\]
The first branch of $F$ depends only on $x$, but both branches of $TFT^{-1}$ depend on both mixed coordinates.
\end{proposition}

\begin{proof}
Since
\[
T^{-1}(u,v)=\left(\frac{u+v}{2},\frac{u-v}{2}\right),
\]
one obtains
\[
TFT^{-1}(u,v)
=
\left(
\frac{(u+v)^2}{4}+\frac{u-v}{2},
\frac{(u+v)^2}{4}-\frac{u-v}{2}
\right).
\]
Both components depend on $u$ and $v$.
\end{proof}

Functional equivalence under an invertible total-state recoding does not preserve the receiver marks or interfaces. Architectural locality belongs to the marked presentation, not to the unmarked function class.
